\documentclass[11pt]{article}
\usepackage[review]{acl}
\usepackage{times}
\usepackage{latexsym}
\usepackage{booktabs}
\usepackage[T1]{fontenc}
\usepackage[utf8]{inputenc}

\title{Section 4 Draft: Experimental Setup}
\author{Anonymous ARR Submission}

\begin{document}
\maketitle

\section{Experimental Setup}
\label{sec:experimental-setup}

We evaluate THEMIS on SWE-bench Lite, a 300-instance subset of real-world GitHub issue-resolution tasks drawn from Python repositories \citep{jimenez-etal-2024-swebench}.  Each instance contains a repository, a base commit, a natural-language problem statement, optional hints, and failing tests used by the official SWE-bench evaluation harness.  THEMIS consumes the issue text and selected source files, produces a patch as a standard \texttt{git diff}, and writes predictions in the JSONL format expected by the harness.

\subsection{Dataset and Instance Processing}
\label{subsec:dataset-processing}

The primary benchmark run uses all 300 SWE-bench Lite test instances with random seed 42.  The instances span the benchmark's Python repository set, with the largest groups coming from Django, SymPy, scikit-learn, Matplotlib, Sphinx, Astropy, Xarray, Flask, Seaborn, Requests, Pytest, and other repositories.  For each instance, the runner checks out the repository at the recorded base commit, selects target files, builds a requirement string, executes the integrated THEMIS workflow, applies the resulting file updates, and emits a patch.

The requirement string concatenates the problem statement, available hints, failing test identifiers, and---for selected semantic-contract presets---additional repair contracts.  This construction is intentionally explicit: failing tests and semantic contracts are not hidden evaluator signals, but become part of the textual requirement object that the graph manager can decompose into requirement nodes.

\subsection{Model and Workflow Configuration}
\label{subsec:model-workflow-config}

All reported runs in this paper use \texttt{gpt-5.1-codex-mini} as both the repair workflow model and the advanced-analysis model.  The main benchmark uses the integrated workflow described in Section~\ref{sec:system-architecture}: Advanced Analysis, knowledge-graph construction, Developer revision, graph rebuilding, and Judge validation.  The default strategy is \texttt{auto\_select}; most controlled variants use \texttt{graph\_only} so that the comparison focuses on structural, context, or semantic-contract changes rather than a changing strategy selector.

The maximum revision count is one in the documented experiments.  Within a revision, the Developer may attempt bounded edit modes, but the workflow-level loop is intentionally shallow.  This setting makes the evaluation closer to a constrained repair pass than to an open-ended autonomous debugging session.

\subsection{File Selection and Context Construction}
\label{subsec:file-selection}

THEMIS uses conservative file selection because the current graph manager is optimized for bounded Python modules.  The runner first extracts explicit Python paths from the issue text and failing-test metadata.  If no direct path is available, it extracts up to a small set of meaningful tokens and searches the repository using a priority order of \texttt{ripgrep}, \texttt{git grep}, and a Python fallback scan.  Candidate files are ranked by token matches, with special weighting for identifier-like tokens, and test files are filtered when selecting the primary repair target.

Several presets expand this conservative target in controlled ways.  Fault-space fallback and neighborhood variants add bounded nearby files; context variants add hand-curated auxiliary files for selected instances; retrieval variants add Python files mentioned in the instance text; and semantic-contract variants append instance-specific repair constraints to the requirement string.  These variants allow us to probe whether additional fault-space or semantic information helps, but they do not remove the bounded-file assumption.

\subsection{Experiment Presets}
\label{subsec:presets}

Table~\ref{tab:presets} summarizes the presets used in the documented experiments.  The full 300-instance benchmark uses the default integrated configuration.  Additional ablation-labeled runs are performed on a fixed five-instance slice.  Because several ablation labels are documented by name rather than by a fully specified component removal, we treat them as probes of selected design choices rather than as a complete causal decomposition of every system component.

\begin{table*}[t]
\centering
\small
\begin{tabular}{llll}
\toprule
Preset family & Workflow / strategy & Scope & Purpose \\
\midrule
Default & integrated, \texttt{auto\_select} & 300 instances & Main benchmark run \\
\texttt{ablation1} & integrated variant, \texttt{graph\_only} & fixed 5-instance slice & graph-only conflict baseline \\
Other ablation labels & ablation-labeled variants & fixed 5-instance slice & probe design choices; definitions require cautious interpretation \\
Fault-space variants & fallback / neighborhood / context / retrieval & selected slices & test bounded target expansion and context additions \\
Semantic contract & contract / control / rerank & controlled semantic slices & test explicit semantic constraints and reranking \\
\bottomrule
\end{tabular}
\caption{Experiment preset families.  We distinguish the full 300-instance benchmark from smaller ablation and semantic-contract slices.}
\label{tab:presets}
\end{table*}

\subsection{Metrics}
\label{subsec:metrics}

We report six metrics.  \textbf{Resolved rate} is the official SWE-bench harness outcome: the number of submitted instances whose generated patch passes the relevant benchmark tests divided by the number of submitted instances.  We also report a completed-only denominator, dividing resolved instances by completed non-empty-patch instances, because empty patches are not meaningful repair attempts.  \textbf{Patch generation rate} is the fraction of submitted instances for which the workflow produces a non-empty patch.  This is an engineering throughput measure, not a correctness measure.  \textbf{Runtime} is wall-clock processing time per instance.  \textbf{Advanced Analysis token usage} is the token count recorded by the separate advanced-analysis interface.  Finally, \textbf{chat-model usage} is not used as reliable cost evidence: the current callback-based accounting records \texttt{chat\_usage.total\_tokens=0}, but this reflects an instrumentation limitation for the chat-model path and should not be interpreted as absence of Developer or Judge LLM calls.

\bibliographystyle{acl_natbib}
\bibliography{../references}

\end{document}
